\FloatBarrier
\section{Esercitazioni di Network Science: Metriche, Assortativit\`a e Modularit\`a}
\label{ch:dm-ex-network-science}

\subsection{Esercizio 1: Caratterizzazione Strutturale di Base}
\label{sec:dm-ex-ns-basic}

Si consideri il grafo non orientato e non pesato di $N = 7$ nodi ($A, \dots, G$) ed $L = 11$ archi ($AB$, $AC$, $BC$, $BD$, $CD$, $CE$, $DE$, $DF$, $EF$, $EG$, $FG$) raffigurato nella \autoref{fig:dm-ex-ns-g1}.

\begin{figure}[H]
\centering
\begin{tikzpicture}[
    v/.style={circle, draw=navy, fill=codebackground, line width=0.9pt, minimum size=6.5mm, font=\small\bfseries},
    e/.style={line width=0.9pt, color=navy!80}
]
    \node[v] (A) at (0,1) {$A$};
    \node[v] (B) at (1.6,2.0) {$B$};
    \node[v] (C) at (1.6,0.0) {$C$};
    \node[v] (D) at (3.4,2.0) {$D$};
    \node[v] (E) at (3.4,0.0) {$E$};
    \node[v] (F) at (5.2,2.0) {$F$};
    \node[v] (G) at (5.2,0.0) {$G$};

    \draw[e] (A)--(B) (A)--(C) (B)--(C) (B)--(D) (C)--(D) (C)--(E) (D)--(E) (D)--(F) (E)--(F) (E)--(G) (F)--(G);
\end{tikzpicture}
\caption{Grafo non orientato di $7$ nodi e $11$ archi per l'Esercizio 1.}
\label{fig:dm-ex-ns-g1}
\end{figure}

\begin{enumerate}
    \item \textbf{Matrice di adiacenza $A$:} ordinando i nodi da $A$ a $G$, la matrice simmetrica $A \in \{0,1\}^{7 \times 7}$ \`e:
    \begin{equation*}
        A = \begin{pmatrix}
        0 & 1 & 1 & 0 & 0 & 0 & 0 \\
        1 & 0 & 1 & 1 & 0 & 0 & 0 \\
        1 & 1 & 0 & 1 & 1 & 0 & 0 \\
        0 & 1 & 1 & 0 & 1 & 1 & 0 \\
        0 & 0 & 1 & 1 & 0 & 1 & 1 \\
        0 & 0 & 0 & 1 & 1 & 0 & 1 \\
        0 & 0 & 0 & 0 & 1 & 1 & 0
        \end{pmatrix}
    \end{equation*}
    \item \textbf{Gradi e grado medio $\langle k \rangle$:} i gradi dei nodi sono $(2, 3, 4, 4, 4, 3, 2)$, con somma $\sum_{i=1}^7 k_i = 22 = 2L$ e grado medio:
    \begin{equation*}
        \langle k \rangle = \frac{2L}{N} = \frac{22}{7} \approx 3.143
    \end{equation*}
    \item \textbf{Densit\`a $\rho$ e distribuzione dei gradi $P(k)$:}
    \begin{equation*}
        \rho = \frac{2L}{N(N-1)} = \frac{22}{7 \times 6} = \frac{11}{21} \approx 0.524
    \end{equation*}
    \begin{align*}
        P(2) &= \frac{2}{7} \text{ ($A,G$)}, \qquad P(3) = \frac{2}{7} \text{ ($B,F$)}, \\
        P(4) &= \frac{3}{7} \text{ ($C,D,E$)}, \qquad P(k) = 0 \text{ altrove.}
    \end{align*}
    \item \textbf{Diametro:} la massima distanza geodetica tra tutte le coppie di nodi \`e pari a \textbf{$3$} (ad esempio lungo il cammino minimo $A - C - E - G$).
\end{enumerate}

% ====================================================================
\FloatBarrier
\subsection{Esercizio 2: Coefficienti di Clustering e Centralit\`a di Grado}
\label{sec:dm-ex-ns-clustering}

Si consideri il grafo non orientato di $N = 7$ nodi formato da due cricche triangolari $\{A, B, C\}$ ed $\{E, F, G\}$ collegate dal nodo ponte $D$ tramite gli archi $CD$ e $DE$ ($8$ archi totali, \autoref{fig:dm-ex-ns-g2}), avente sequenza dei gradi $(2, 2, 3, 2, 3, 2, 2)$.

\begin{figure}[H]
\centering
\begin{tikzpicture}[
    v/.style={circle, draw=navy, fill=codebackground, line width=0.9pt, minimum size=6.5mm, font=\small\bfseries},
    e/.style={line width=0.9pt, color=navy!80}
]
    \node[v] (A) at (0,1.6) {$A$};
    \node[v] (B) at (0,0.0) {$B$};
    \node[v] (C) at (1.5,0.8) {$C$};
    \node[v, draw=crimson] (D) at (3.2,0.8) {$D$};
    \node[v] (E) at (4.9,0.8) {$E$};
    \node[v] (F) at (6.4,1.6) {$F$};
    \node[v] (G) at (6.4,0.0) {$G$};

    \draw[e] (A)--(B) (A)--(C) (B)--(C) (C)--(D) (D)--(E) (E)--(F) (E)--(G) (F)--(G);
\end{tikzpicture}
\caption{Grafo formato da due cricche triangolari collegate dal nodo ponte $D$ (Esercizio 2).}
\label{fig:dm-ex-ns-g2}
\end{figure}

\begin{enumerate}
    \item \textbf{Coefficiente di clustering locale $C_v$ e medio $\overline{C}$:}
    Per un nodo $v$ di grado $k_v \ge 2$, $C_v = \frac{2 n_e(v)}{k_v(k_v - 1)}$, dove $n_e(v)$ \`e il numero di archi tra i vicini di $v$:
    \begin{itemize}
        \item $C_A = C_B = C_F = C_G = \frac{2 \cdot 1}{2 \cdot 1} = 1$;
        \item $C_C = \frac{2 \cdot 1}{3 \cdot 2} = \frac{1}{3}$ (vicini $\{A,B,D\}$, solo $AB$ \`e presente);
        \item $C_D = \frac{2 \cdot 0}{2 \cdot 1} = 0$ (vicini $\{C,E\}$ non collegati);
        \item $C_E = \frac{2 \cdot 1}{3 \cdot 2} = \frac{1}{3}$ (vicini $\{D,F,G\}$, solo $FG$ \`e presente).
    \end{itemize}
    Il coefficiente di clustering medio \`e quindi:
    \begin{equation*}
        \overline{C} = \frac{1 + 1 + 1 + 1 + \frac{1}{3} + 0 + \frac{1}{3}}{7} = \frac{14/3}{7} = \frac{2}{3}
    \end{equation*}
    \item \textbf{Coefficiente di clustering globale $C_{\text{global}}$:} il numero totale di triplette connesse \`e:
    \begin{equation*}
        N_{\text{triplets}} = \sum_{v \in V} \binom{k_v}{2} = \frac{1}{2}\left[4 \cdot (2 \cdot 1) + 2 \cdot (3 \cdot 2) + 1 \cdot (2 \cdot 1)\right] = \frac{22}{2} = 11
    \end{equation*}
    Poich\'e il grafo contiene $N_{\text{triangles}} = 2$ triangoli ($ABC$ ed $EFG$), si hanno $N_{\text{closed triplets}} = 3 \times 2 = 6$, da cui:
    \begin{equation*}
        C_{\text{global}} = \frac{3 N_{\text{triangles}}}{N_{\text{triplets}}} = \frac{6}{11}
    \end{equation*}
    \item \textbf{Centralit\`a di grado normalizzata:} definita come $C_D^{\text{deg}}(v) = \frac{k_v}{N - 1}$. Con $N - 1 = 6$, per $C$ e $D$ si ottiene:
    \begin{equation*}
        C_D^{\text{deg}}(C) = \frac{3}{6} = \frac{1}{2}, \qquad C_D^{\text{deg}}(D) = \frac{2}{6} = \frac{1}{3}
    \end{equation*}
\end{enumerate}

% ====================================================================
\FloatBarrier
\subsection{Esercizio 3: Assortativit\`a Categoriale di Newman}
\label{sec:dm-ex-ns-assortativity}

Si consideri il grafo non orientato di $8$ nodi della \autoref{fig:dm-ex-ns-g3}, ripartito in $4$ nodi verdi ($G_1, \dots, G_4$) e $4$ nodi rossi ($R_1, \dots, R_4$).

\begin{figure}[H]
\centering
\begin{tikzpicture}[
    vg/.style={circle, draw=darkgreen, fill=darkgreen!15, line width=0.9pt, minimum size=6.5mm, font=\scriptsize\bfseries},
    vr/.style={circle, draw=crimson, fill=crimson!15, line width=0.9pt, minimum size=6.5mm, font=\scriptsize\bfseries},
    e/.style={line width=0.9pt, color=navy!75}
]
    \node[vg] (G1) at (0,1.6) {$G_1$};
    \node[vg] (G2) at (0,0.0) {$G_2$};
    \node[vg] (G3) at (1.8,1.6) {$G_3$};
    \node[vg] (G4) at (1.8,0.0) {$G_4$};

    \node[vr] (R1) at (4.0,1.6) {$R_1$};
    \node[vr] (R2) at (5.8,1.6) {$R_2$};
    \node[vr] (R3) at (4.0,0.0) {$R_3$};
    \node[vr] (R4) at (5.8,0.0) {$R_4$};

    \draw[e] (G1)--(G2) (G1)--(G3) (G2)--(G3) (G2)--(G4) (G3)--(G4);
    \draw[e] (R1)--(R2) (R1)--(R3) (R2)--(R3) (R2)--(R4) (R3)--(R4);
    \draw[e, dashed, navy] (G3)--(R1) (G4)--(R3);
\end{tikzpicture}
\caption{Rete bicolore con $L_{GG}=5$ archi verde-verde, $L_{RR}=5$ rosso-rosso e $L_{GR}=2$ misti (Esercizio 3).}
\label{fig:dm-ex-ns-g3}
\end{figure}

Per un grafo non orientato ($a_i = b_i$), il coefficiente di assortativit\`a categoriale di Newman \`e:
\begin{equation}
    r = \frac{\sum_i e_{ii} - \sum_i a_i^2}{1 - \sum_i a_i^2}
\end{equation}
Svolgendo i quattro passi di calcolo:
\begin{enumerate}
    \item \textbf{Conteggio degli archi:} $L_{GG} = 5$, $L_{RR} = 5$, $L_{GR} = 2$, per un totale di $L = 12$ archi ($2L = 24$ estremit\`a).
    \item \textbf{Frazioni intra-categoria:} $e_{GG} = \frac{5}{12}$, $e_{RR} = \frac{5}{12} \implies \sum_i e_{ii} = \frac{10}{12}$.
    \item \textbf{Frazioni di estremit\`a:}
    \begin{align*}
        a_G &= \frac{2 L_{GG} + L_{GR}}{2L} = \frac{10 + 2}{24} = \frac{1}{2}, \\
        a_R &= \frac{2 L_{RR} + L_{GR}}{2L} = \frac{10 + 2}{24} = \frac{1}{2}
    \end{align*}
    \item \textbf{Calcolo di $r$:}
    \begin{equation*}
        r = \frac{\frac{10}{12} - \left[\left(\frac{1}{2}\right)^2 + \left(\frac{1}{2}\right)^2\right]}{1 - \left[\left(\frac{1}{2}\right)^2 + \left(\frac{1}{2}\right)^2\right]} = \frac{\frac{10}{12} - \frac{1}{2}}{1 - \frac{1}{2}} = \frac{\frac{1}{3}}{\frac{1}{2}} = \frac{2}{3} \approx 0.667
    \end{equation*}
\end{enumerate}
Il valore positivo $r \approx 0.667$ indica una forte \textbf{assortativit\`a categoriale positiva} (omofilia): i nodi tendono a connettersi preferenzialmente con nodi della stessa categoria rispetto al modello nullo casuale.

% ====================================================================
\FloatBarrier
\subsection{Esercizio 4: Modularit\`a e Partizioni in Comunit\`a}
\label{sec:dm-ex-ns-modularity}

Si consideri il grafo non orientato di $N = 8$ nodi ($A, \dots, H$) ed $L = 12$ archi ($2L = 24$) con gradi $(3, 2, 4, 3, 4, 3, 3, 2)$ della \autoref{fig:dm-ex-ns-g4}, sul quale si confrontano due partizioni: $\mathcal{P}_1 = \{\{A,B,C,D\}, \{E,F,G,H\}\}$ e $\mathcal{P}_2 = \{\{A,B,C\}, \{D,E,F,G,H\}\}$.

\begin{figure}[H]
\centering
\begin{tikzpicture}[
    v/.style={circle, draw=navy, fill=codebackground, line width=0.9pt, minimum size=6mm, font=\scriptsize\bfseries},
    e/.style={line width=0.9pt, color=navy!75}
]
    \node[v] (A) at (0,1.6) {$A$};
    \node[v] (B) at (0,0.0) {$B$};
    \node[v] (C) at (1.5,1.6) {$C$};
    \node[v] (D) at (1.5,0.0) {$D$};
    \node[v] (E) at (3.6,1.6) {$E$};
    \node[v] (F) at (3.6,0.0) {$F$};
    \node[v] (G) at (5.1,1.6) {$G$};
    \node[v] (H) at (5.1,0.0) {$H$};

    \draw[e] (A)--(B) (A)--(C) (A)--(D) (B)--(C) (C)--(D);
    \draw[e] (E)--(F) (E)--(G) (E)--(H) (F)--(G) (G)--(H);
    \draw[e] (C)--(E) (D)--(F);

    \draw[dashed, darkgreen, rounded corners=5pt, line width=0.9pt] (-0.5,-0.5) rectangle (2.0,2.1);
    \draw[dashed, darkgreen, rounded corners=5pt, line width=0.9pt] (3.1,-0.5) rectangle (5.6,2.1);
    \node[font=\scriptsize\bfseries, color=darkgreen] at (2.55,-0.85) {Partizione $\mathcal{P}_1$: $L_{c_1} = 5$, $L_{c_2} = 5$};
\end{tikzpicture}
\caption{Grafo di $8$ nodi e $12$ archi con evidenziata la partizione $\mathcal{P}_1$ (Esercizio 4).}
\label{fig:dm-ex-ns-g4}
\end{figure}

Raggruppando per comunit\`a $c$ i termini della definizione generale di modularit\`a $Q = \frac{1}{2L}\sum_{i,j}\bigl(A_{ij} - \frac{k_i k_j}{2L}\bigr)\delta(c_i, c_j)$, poich\'e $\sum_{i,j \in c} A_{ij} = 2L_c$ e $\sum_{i,j \in c} k_i k_j = \bigl(\sum_{i \in c} k_i\bigr)^2 = d_c^2$, si ottiene la formulazione a livello di comunit\`a:
\begin{equation}
    Q = \sum_{c} \left[ \frac{L_c}{L} - \left(\frac{d_c}{2L}\right)^2 \right]
\end{equation}
dove $L_c$ \`e il numero di archi interni alla comunit\`a $c$ e $d_c = \sum_{i \in c} k_i$ \`e la somma dei gradi dei nodi in $c$.

\begin{itemize}
    \item \textbf{Partizione $\mathcal{P}_1$:} entrambe le comunit\`a $\{A,B,C,D\}$ ed $\{E,F,G,H\}$ hanno $L_c = 5$ e $d_c = 12$:
    \begin{equation*}
        Q(\mathcal{P}_1) = 2 \left[ \frac{5}{12} - \left(\frac{12}{24}\right)^2 \right] = 2 \left[ \frac{5}{12} - \frac{1}{4} \right] = \frac{1}{3} \approx 0.333
    \end{equation*}
    \item \textbf{Partizione $\mathcal{P}_2$:} $\{A,B,C\}$ ha $L_c = 3, d_c = 9$, mentre $\{D,E,F,G,H\}$ ha $L_c = 6, d_c = 15$:
    \begin{equation*}
        Q(\mathcal{P}_2) = \left[ \frac{3}{12} - \left(\frac{9}{24}\right)^2 \right] + \left[ \frac{6}{12} - \left(\frac{15}{24}\right)^2 \right] = \frac{7}{32} = 0.21875
    \end{equation*}
\end{itemize}

Poich\'e $Q(\mathcal{P}_1) = \frac{1}{3} > Q(\mathcal{P}_2) = \frac{7}{32}$, \textbf{$\mathcal{P}_1$ \`e la partizione migliore}. Il solo conteggio degli archi interni non \`e sufficiente a giudicare una partizione, perch\'e la modularit\`a sottrae la connettivit\`a attesa sotto un modello nullo che preserva i gradi. Algoritmi di \textit{community detection} come \textbf{Louvain} mirano appunto a massimizzare $Q$ spostando iterativamente i nodi tra comunit\`a adiacenti e aggregando le comunit\`a in super-nodi.
